\begin{statement}{D.26}{Proved}{The exact covering restriction of the central peripheral obstruction}
Keep the finite groups, epimorphism $q_G:\widehat F\to G$,
factor-swap automorphism $\alpha$, and central extension
$\pi:E\to G$ of D.24.  Put $\theta=\alpha q_G$ and
$S_i=\overline{\langle x_i\rangle}$, with
$x_3=(x_1x_2)^{-1}$.  For any open subgroup
$H\leq\widehat F$ of index $d$, take representatives
$t\in F=\langle x_1,x_2\rangle$ of
$H\backslash\widehat F/S_i$ and set
\[
 e_{i,t}=[S_i:S_i\cap t^{-1}Ht],\qquad
 h_{i,t}=t x_i^{e_{i,t}}t^{-1},\qquad
 S_{i,t}=\overline{\langle h_{i,t}\rangle}\leq H.
 \tag{D.26.1}
\]
Let $b_H$ be the number of these induced peripheral subgroups.
The number of continuous homomorphisms $v:H\to E$ satisfying
\[
       \pi v=\theta|_H,\qquad
       v(h_{i,t})\sim_E X_i^{e_{i,t}}
                  \quad\text{for every }(i,t)
 \tag{D.26.2}
\]
is exactly
\[
 \begin{cases}
  0,&d\text{ odd},\\
  2^{\,d+2-b_H},&d\text{ even}.
 \end{cases}
 \tag{D.26.3}
\]
These are literal weak lifts with the fixed base map; no
surjectivity or quotient by conjugation is imposed on the count.

In particular every cover of odd index retains the nonzero
peripheral obstruction of D.24.  For every odd prime $\ell$,
this includes every open subgroup obtained as the inverse image
of a subgroup of a finite $\ell$-group quotient of $\widehat F$.
The exact map controlling this assertion is restriction
\[
\begin{split}
 H^2_{\cts}(\widehat F,\{S_i\};\F_2)&\longrightarrow
           H^2_{\cts}(H,\{S_{i,t}\};\F_2),\\
                   \F_2\ni a&\longmapsto d a\in\F_2,
\end{split}
 \tag{D.26.4}
\]
where both identifications use the sum of the oriented boundary
coordinates.  Thus passage to deeper covers of odd index cannot
solve this specified finite peripheral problem.
\end{statement}
\begin{proof}
The dense discrete free group $F$ acts transitively on the
$d$ cosets of $H$.  Its subgroup $H_0=F\cap H$ has index $d$,
and its profinite completion maps isomorphically onto $H$.
Indeed the core in $F$ of every finite-index subgroup of $H_0$
has finite index in $F$, proving that the induced topology is the
full profinite topology.  The discrete groups
$\langle h_{i,t}\rangle$ complete to precisely the groups in
\textup{(D.26.1)}.  The indices $e_{i,t}$ are the orbit lengths
of $S_i$ on the $d$ cosets, so
\[
                  \sum_t e_{i,t}=d\quad(1\leq i\leq3).
 \tag{D.26.5}
\]

Realize $F$ as the fundamental group of an oriented sphere
with three open discs removed, taking the oriented boundary
loops to be $x_1,x_2,x_3$.  Its connected cover associated to
$H_0$ is an oriented compact surface with $b_H$ boundary
components.  These components give exactly the loops
$h_{i,t}$ with their specified orientations and base paths.
Cutting this surface along a system of handle and boundary
arcs gives the presentation
\[
 \left\langle a_1,b_1,\ldots,a_g,b_g,
      c_1,\ldots,c_{b_H}\ \middle|\
      \prod_{j=1}^{g}[a_j,b_j]\prod_{k=1}^{b_H}c_k=1
 \right\rangle,
 \tag{D.26.6}
\]
where the $c_k$ are conjugates of the displayed induced boundary
generators.  Abelianization therefore gives exactly one relation
among their classes, namely $\sum_k c_k=0$.  Since
$H_0$ is free of rank $d+1$, comparison of ranks yields
$2g=d+2-b_H$.

For trivial coefficients $\F_2$, continuous first cohomology of
the free profinite and procyclic groups equals the first
cohomology of their discrete dense free and cyclic groups:
both are homomorphisms to the same finite group.
Their second cohomology vanishes.  The relative long exact
sequences consequently identify both relative second groups
with the cokernels of their boundary restriction maps.
By \textup{(D.26.6)}, the latter map for $H$ has image
\[
       \{(u_k)\in\F_2^{b_H}:\textstyle\sum_k u_k=0\},
 \tag{D.26.7}
\]
and kernel $\F_2^{2g}$.  For the base pair the corresponding
map is $(a,b)\mapsto(a,b,-a-b)$.
These calculations also prove the discrete-to-profinite
comparison for the relative groups being used.

Under restriction, a boundary coordinate $u_i$ pulls back to
$e_{i,t}u_i$ on $S_{i,t}$; conjugation of base paths has no
effect on these trivial coefficients.  Summing the coordinates
and using \textup{(D.26.5)} gives
$\sum_{i,t}e_{i,t}u_i=d\sum_i u_i$.  Naturality of the relative
long exact sequence proves exactly \textup{(D.26.4)}, including
its source, target and the chosen identifications.

The relative extension class of D.24, with local lifts
$U_i=(B_i,a_i)$, is $1\in\F_2$ because
$U_1U_2U_3=(-I,1)$.  For a based induced peripheral group,
choose any lift $\Theta_t\in E$ of $\theta(t)$ and use
$\Theta_t U_i^{e_{i,t}}\Theta_t^{-1}$ as its local image.
Changing $\Theta_t$ by the central kernel changes nothing.
Its projection is $\theta(h_{i,t})$, and it is conjugate to
$X_i^{e_{i,t}}$ by D.24.  These are precisely the restricted
local maps in the relative pullback extension, so their relative
obstruction is $d\in\F_2$.

We check the full conjugacy freedom in \textup{(D.26.2)}.
Every $X_i$ has odd order.  In a central extension by $C_2$,
the two elements above any odd-order base element differ by the
central involution, and exactly one has odd order.  Thus the
displayed local image is the unique odd-order lift of
$\theta(h_{i,t})$.  Every conjugate allowed in
\textup{(D.26.2)} has odd order and equals that lift.
Equivalently, each induced base-centralizer commutator is killed
both by two and by an odd order, hence is trivial.  This remains
true for conjugators in all of $E$ and for a proper image
$\theta(H)\leq G$.  The relative obstruction therefore retains
the complete unframed peripheral condition.

For completeness, existence and the count follow directly from
cochains.  The pullback of $E$ to the free profinite group $H$
has a global section: lift its $d+1$ free generators arbitrarily.
The discrepancies between this section and the prescribed local
maps form a vector in $\F_2^{b_H}$.  Multiplication of the
section by a character $H\to C_2$ changes that vector by the
boundary restriction map.  Its sum is the relative class $d$.
Equation \textup{(D.26.7)} gives no correction for odd $d$ and
a torsor under $\F_2^{2g}$ for even $d$.  This proves the exact
count \textup{(D.26.3)} without imposing an extra generation
condition.

The program \texttt{d\_peripheral\_covers.py} constructs
connected coset covers of indices one through nine, with cyclic
and reflection permutation data.  For each it computes literal
Schreier basis words and every induced boundary word.  It
enumerates every central correction to the lifted basis and
substitutes every surviving lift into all boundary equations
in $\operatorname{SL}_2(\F_{11})$.  The computed counts agree
with \textup{(D.26.3)}, including positive-genus even covers.
The code, full coset and word data, every correction, and the
executed output are retained.  The all-index conclusion follows
from the restriction calculation above.
\end{proof}
